\documentclass[11pt,a4paper]{article}
\usepackage[margin=25mm]{geometry}
\usepackage{amsmath,amssymb,amsthm,booktabs}
\usepackage[T1]{fontenc}
\usepackage{lmodern}
\usepackage[hidelinks]{hyperref}
\setlength{\emergencystretch}{3em}
\newtheorem{proposition}{Proposition}
\title{Fine numbers are not log-concave\\\large A counterexample to Conjecture 00000003327}
\author{ziangni-sys}
\date{4 October 2026}

\begin{document}
\maketitle
\begin{abstract}
The positive consecutive Fine numbers $F_2=1$, $F_3=2$, $F_4=6$
violate log-concavity: $F_3^2=4<6=F_2F_4$.
We derive these values from hill-free Dyck paths and supply a
self-contained Lean 4 verification of the counting definitions and
the strict inequality.
\end{abstract}

\section{The claim and its interpretation}
Conjecture 00000003327 asserts log-concavity of the Fine number
sequence, in the context of group gradings and a proposed Narayana
convolution explanation. We use the standard Fine numbers indexed
from $F_0=1$: $F_n$ counts hill-free Dyck paths of semilength $n$.
This interpretation is stated explicitly in the introduction of
Barcucci et al.~\cite{hillfree}.

A Dyck path is a word in $U=(1,1)$ and $D=(1,-1)$ starting at height
zero, never attaining negative height, and ending at height zero.
A \emph{hill} is a consecutive $UD$ starting at height zero, thus
forming a peak of height one. A path is hill-free if it contains
no such pair.

Log-concavity of a nonnegative sequence requires
\[
  F_n^2\geq F_{n-1}F_{n+1}
\]
at each interior index. We disprove even the weaker requirement
restricted to $n\geq3$, entirely within the positive tail. The zero
$F_1$ plays no role. Shifting the indexing convention also leaves
the consecutive triple $1,2,6$ and its violation unchanged.
The conjecture specifies no alternative normalization, no different
sequence of graded dimensions, and no eventual-only qualification.
Our conclusion concerns its stated Fine-number log-concavity clause;
falsifying that clause is sufficient, independently of the suggested
proof mechanism.

\section{Exact counterexample}
\begin{proposition}
The Fine number sequence is not log-concave, even on its positive tail.
\end{proposition}
\begin{proof}
The complete hill-free lists for semilengths $2$, $3$, and $4$ are:
\begin{center}
\begin{tabular}{clc}
\toprule
$n$ & Hill-free Dyck paths & $F_n$ \\
\midrule
2 & \texttt{UUDD} & 1 \\
3 & \texttt{UUUDDD}, \texttt{UUDUDD} & 2 \\
4 & \texttt{UUUUDDDD}, \texttt{UUUDUDDD}, \texttt{UUUDDUDD} & \\
  & \texttt{UUDUUDDD}, \texttt{UUDUDUDD}, \texttt{UUDDUUDD} & 6 \\
\bottomrule
\end{tabular}
\end{center}
To see that the lists are exhaustive, split a path at each return
to height zero. A hill-free irreducible component has the unique form
$UAD$, where $A$ is a nonempty Dyck path. For semilengths $2$ and $3$
there can only be one component, giving respectively the one and two
Dyck interiors of semilength $1$ and $2$. For semilength $4$, a single
component gives the five Dyck interiors of semilength $3$; two
components must both have semilength $2$, giving one further path.
Therefore $F_2=1$, $F_3=2$, $F_4=6$, and
$F_3^2=4<6=F_2F_4$.
\end{proof}

\section{Formal verification and reproducibility}
The accompanying \texttt{lean4/} project uses Lean 4.19.0 and its
standard library only. In \texttt{Main.lean}, Boolean lists encode
step words, with \texttt{true} for $U$ and \texttt{false} for $D$.
\begin{itemize}
\item \texttt{words} enumerates every word of a prescribed length.
Theorems \texttt{mem\_words} and \texttt{words\_nodup} prove
completeness and absence of repetitions for every length.
\item \texttt{DyckFrom} recursively states that the path stays
nonnegative and ends at zero. \texttt{NoHillsFrom} forbids a $UD$
pair at height zero. The equivalence \texttt{hillFreeFrom\_iff}
verifies that the Boolean scanner accepts exactly their conjunction.
\item \texttt{paths n} filters words of length $2n$ by that scanner.
\texttt{mem\_paths\_iff} identifies precisely the admissible paths,
and \texttt{paths\_nodup} guarantees that length counts distinct paths.
Thus \texttt{fine n := (paths n).length} implements the definition of $F_n$.
\item \texttt{fine\_two}, \texttt{fine\_three}, and \texttt{fine\_four}
calculate the required values by kernel-checked \texttt{decide}.
\texttt{violation} proves the strict inequality, and
\texttt{not\_logConcavePositiveTail} negates the universal claim
for $n\geq3$. Finally, \texttt{conjecture\_00000003327\_false}
negates the full log-concavity clause for $n\geq1$.
\end{itemize}

The printed axiom audit reports no axioms for the three numerical
equalities, the strict inequality, or the final negation.
The general enumeration/semantics helpers use only Lean's standard
logical axioms (\texttt{propext}, \texttt{Quot.sound}, and, for the
semantics equivalence, \texttt{Classical.choice}). There are no added
axioms, admitted proofs, or native-evaluation proof shortcuts.

From the submission directory, run:
\begin{verbatim}
cd lean4
lake build
lake env lean Main.lean
cd ..
python verify.py
pdflatex -interaction=nonstopmode -halt-on-error report.tex
pdflatex -interaction=nonstopmode -halt-on-error report.tex
\end{verbatim}
Alternatively, \texttt{tectonic report.tex} builds the PDF.
The independent Python script uses signed prefix heights to enumerate
all binary step words. It also checks the first-return recurrence
\[
  F_0=1,\qquad
  F_n=\sum_{k=2}^{n}C_{k-1}F_{n-k}\quad(n\geq1),
  \qquad C_j=\frac{1}{j+1}\binom{2j}{j},
\]
against the enumeration for $0\leq n\leq7$.
Python output is supplementary; the Lean contradiction does not
depend on it or on imported numerical data.

\begin{thebibliography}{1}
\bibitem{hillfree}
E. Barcucci, E. Pergola, R. Pinzani, and S. Rinaldi,
\emph{ECO method and hill-free generalized Motzkin paths},
S\'eminaire Lotharingien de Combinatoire \textbf{46} (2001),
Article B46b, p.~1.
\par\url{https://www.mat.univie.ac.at/~slc/wpapers/s46rinaldi.pdf}.
\end{thebibliography}
\end{document}
